\documentclass[12pt, a4paper]{article}

\usepackage[utf8]{inputenc}
\usepackage[T1]{fontenc}
\usepackage{amsmath, amssymb, amsthm}
\usepackage{geometry}
\usepackage{fancyhdr}
\usepackage{enumitem}
\usepackage{xcolor}
\usepackage{hyperref}

\geometry{margin=1in}
\pagestyle{fancy}
\fancyhf{}
\rhead{Mathematics}
\lhead{Student}
\rfoot{Page \thepage}

\newcommand{\problem}[1]{\subsection*{Problem #1}}
\newcommand{\solution}{\paragraph{Solution.}}

\definecolor{solutioncolor}{RGB}{0, 100, 0}

\begin{document}

\begin{center}
    {\LARGE\bfseries Homework Solutions} \\[0.5em]
    {\large Mathematics} \\[0.3em]
    Student \quad | \quad 2026-10-08
\end{center}

\hrule
\vspace{1em}
\problem{Q1}

Suppose $f$ is a function and $L$ and $a$ are real numbers.

\solution

See sub-problems for individual limit computations.

\textbf{Answer:}
\[\boxed{\text{See sub-problems}}\]

\textbf{(a)}
$\lim_{x \to a} f(x) = L$ means: for every $\varepsilon > 0$ there exists $\delta > 0$ such that
$0 < |x - a| < \delta \Rightarrow |f(x) - L| < \varepsilon$.
Informally: we can make $f(x)$ as close to $L$ as we like by taking $x$ sufficiently close to (but not equal to) $a$.
$\boxed{\forall \varepsilon > 0, \exists \delta > 0: 0 < |x-a| < \delta \Rightarrow |f(x)-L| < \varepsilon}$

\textbf{(b)}
This requires graphing, which is a visual exercise.
$\boxed{\text{See graph}}$

\textbf{(c)}
This requires graphing, which is a visual exercise.
$\boxed{\text{See graph}}$

\textbf{(d)}
This requires graphing, which is a visual exercise.
$\boxed{\text{See graph}}$

\vspace{1em}
\hrule
\vspace{1em}

\problem{Q2}

Is $\infty$ a number? What does $\lim_{x \to a} f(x) = \infty$ mean? Does $\lim_{x \to a} f(x) = \infty$ mean that the limit exists?

\solution

$\infty$ is \textbf{not} a real number.

$\lim_{x \to a} f(x) = \infty$ means that $f(x)$ grows without bound as $x \to a$.

The limit does \textbf{not exist} in the usual sense; we say the limit ``equals infinity'' as a shorthand for divergence.

\textbf{Answer:}
\[\boxed{\infty \text{ is not a number. The limit does not exist.}}\]

\vspace{1em}
\hrule
\vspace{1em}

\problem{Q3}

What does the Squeeze Theorem say? Draw a picture illustrating this theorem.

\solution

\textbf{Squeeze Theorem:} If $g(x) \le f(x) \le h(x)$ for all $x$ near $a$ (except possibly at $a$), and $\lim_{x \to a} g(x) = \lim_{x \to a} h(x) = L$, then $\lim_{x \to a} f(x) = L$.

Intuitively: if $f$ is ``squeezed'' between two functions that both converge to $L$, then $f$ must also converge to $L$.

\textbf{Answer:}
\[\boxed{\text{If } g(x) \le f(x) \le h(x) \text{ and } \lim g = \lim h = L, \text{ then } \lim f = L.}\]

\vspace{1em}
\hrule
\vspace{1em}

\problem{Q4}

What do we mean by $\lim_{x \to a^-} f(x) = L$? What is meant by $\lim_{x \to a^+} f(x) = L$? Graph a function $y = f(x)$ and some point $x = a$ where $\lim_{x \to a^-} f(x) \neq \lim_{x \to a^+} f(x)$. Find an algebraic expression for such a function.

\solution

$\lim_{x \to a^-} f(x) = L$: the \textbf{left-hand limit}. As $x$ approaches $a$ from values \emph{less than} $a$, $f(x)$ approaches $L$.

$\lim_{x \to a^+} f(x) = L$: the \textbf{right-hand limit}. As $x$ approaches $a$ from values \emph{greater than} $a$, $f(x)$ approaches $L$.

The two-sided limit $\lim_{x \to a} f(x) = L$ exists if and only if both one-sided limits exist and are equal.

Example: $f(x) = \begin{cases} x & x < 0 \\ x+1 & x \ge 0 \end{cases}$ at $a = 0$: $\lim_{x \to 0^-} f(x) = 0 \neq 1 = \lim_{x \to 0^+} f(x)$.

\textbf{Answer:}
\[\boxed{\lim_{x \to a^-} f(x): \text{from below}; \quad \lim_{x \to a^+} f(x): \text{from above}}\]

\vspace{1em}
\hrule
\vspace{1em}

\problem{P1}

(a) Is the limit of a sum always the sum of the limits? Give an example where $\lim_{x \to a}[f(x) + g(x)]$ exists even though neither $\lim_{x \to a} f(x)$ nor $\lim_{x \to a} g(x)$ exist.

\solution

No, the limit of a sum is \emph{not} always the sum of the limits.

The \textbf{Sum Law} requires both individual limits to exist: $\lim[f(x)+g(x)] = \lim f(x) + \lim g(x)$ \emph{if both limits exist}.

\textbf{Counterexample:} Let $f(x) = \frac{1}{x}$ and $g(x) = -\frac{1}{x}$. Neither $\lim_{x \to 0} f(x)$ nor $\lim_{x \to 0} g(x)$ exists, but $\lim_{x \to 0}[f(x)+g(x)] = \lim_{x \to 0} 0 = 0$.

\textbf{Answer:}
\[\boxed{\text{No. } f(x)=1/x,\; g(x)=-1/x: \lim[f+g]=0 \text{ but } \lim f, \lim g \text{ DNE}}\]

\vspace{1em}
\hrule
\vspace{1em}

\problem{P2}

(a) Find $\lim_{x \to 0^+}(\frac{1}{x} - \frac{1}{|x|})$, $\lim_{x \to 0^-}(\frac{1}{x} - \frac{1}{|x|})$, and $\lim_{x \to 0}(\frac{1}{x} - \frac{1}{|x|})$, if they exist.

\solution

$\lim_{x \to 0^+} - \frac{1}{\left|{x}\right|} + \frac{1}{x} = 0$

$\lim_{x \to 0^-} - \frac{1}{\left|{x}\right|} + \frac{1}{x} = -\infty$

$\lim_{x \to 0} - \frac{1}{\left|{x}\right|} + \frac{1}{x} = 0$

\textbf{Answer:}
\[\boxed{0}\]

\textbf{(b)}
Computing: $\lim_{x \to 0} \left|{x}\right| = 0$
This confirms the statement. \qed
$\boxed{0}$

\textbf{(c)}
Left limit: $\lim_{x \to 0^-} \frac{\left|{x}\right|}{x} = -1$
Right limit: $\lim_{x \to 0^+} \frac{\left|{x}\right|}{x} = 1$
Since $-1 \neq 1$, the limit does not exist.
$\boxed{\text{DNE}}$

\vspace{1em}
\hrule
\vspace{1em}

\problem{P3}

For what kinds of functions can one evaluate the limit $\lim_{x \to a} f(x)$ by just plugging in $x = a$? Give an example of a function $f(x)$ for which the limit $\lim_{x \to a} f(x)$ exists and can be evaluated, yet for which this substitution doesn't work.

\solution

\textbf{Continuous functions} allow direct substitution: $\lim_{x \to a} f(x) = f(a)$.

This includes polynomials, rational functions (at non-zero denominator points), trigonometric, exponential, and logarithmic functions at points in their domain.

Example where limit exists but substitution fails: $f(x) = \frac{x^2 - 1}{x - 1}$ at $x = 1$. Direct substitution gives $\frac{0}{0}$, but $\lim_{x \to 1} \frac{x^2-1}{x-1} = \lim_{x \to 1}(x+1) = 2$.

\textbf{Answer:}
\[\boxed{\text{Continuous functions: } \lim_{x \to a} f(x) = f(a)}\]

\vspace{1em}
\hrule
\vspace{1em}

\problem{AP1}

Below are listed three limits. Match the best method to each limit and evaluate each limit.

\solution

See sub-problems for individual limit computations.

\textbf{Answer:}
\[\boxed{\text{See sub-problems}}\]

\textbf{(a)}
$\lim_{x \to 0} \frac{x^{1000} - 703 x^{506} + \pi x^{17} + 480}{- 372 x^{75} + 39 x^{14} + \sqrt{7} x^{5} - 24} = -20$
$\boxed{-20}$

\textbf{(b)}
$\lim_{x \to 4} \frac{x^{2} + 2 x - 24}{x - 4} = 10$
$\boxed{10}$

\textbf{(c)}
$\lim_{x \to 0} x^{2} \sin{\left(\frac{1}{x} \right)} = 0$
$\boxed{0}$

\vspace{1em}
\hrule
\vspace{1em}

\problem{AP2}

Using a calculator, try to guess $\lim_{n \to \infty} n \sin(\frac{\pi}{n})$. (Remember that $n$ is in radians, not degrees!)

\solution

$\lim_{x \to \infty} n \sin{\left(\frac{\pi}{n} \right)} = \pi$

\textbf{Answer:}
\[\boxed{\pi}\]

\vspace{1em}
\hrule
\vspace{1em}

\problem{AP3}

Draw a circle of radius 1. Inscribe a regular $n$-gon inside it.

\textit{\color{red} 未求解: 无法提取极限表达式}

\textbf{(a)}
\textit{\color{red} 无法解析子题极限}

\textbf{(b)}
\textit{\color{red} 无法解析子题极限}

\textbf{(c)}
\textit{\color{red} 无法解析子题极限}

\vspace{1em}
\hrule
\vspace{1em}


\vspace{2em}
\noindent\textit{Auto-generated: 9/10 problems solved by SymPy.}
\end{document}
